% ===================================================================== Chapter 14 (continued)
% Additional nonlinear buckling analysis using ANSYS LS-DYNA (follow-on extension; values from
% 15_LS_DYNA_Extension/12C_Nonlinear_Buckling: FINAL_12C_RESULTS.md, IMPERFECTION_SENSITIVITY.csv, MECHANICAL_vs_LSDYNA.csv)
\section{Additional Nonlinear Buckling Analysis Using ANSYS LS-DYNA}\label{sec:lsdyna}
This section reports an additional nonlinear buckling analysis using ANSYS LS-DYNA, carried out as a follow-on extension of the
thermo-structural analysis of Chapters~\ref{ch:buckling} and~\ref{ch:supports}. It is not part of the internship work programme described in
Chapter~\ref{ch:intro} and does not change any result of the preceding chapters; the Mechanical linear buckling result ($\lambda_1$ = 1.108,
$P_{cr}$ = 608.25~kN under S1) remains the reference.

\subsection{Purpose and Model}
The linear eigenvalue analysis gives the bifurcation load of a perfectly straight, elastic column (Chapter~\ref{ch:buckling}). It cannot show how a duct
with an initial out-of-straightness deflects as the thermal load rises. The extension therefore follows the geometrically nonlinear
(large-deflection) load path of the S1 / LC2 configuration with small imperfections in the shape of the first buckling mode.

The LS-DYNA model is the LC2 model of Chapter~\ref{ch:struct} carried over unchanged:
\begin{itemize}
\item the same mesh (108,252 nodes, 23,400 twenty-node solid elements);
\item the same CFD-derived temperature field mapped from the medium Fluent solution, with the reference temperature of 300~K;
\item the same LC2 support condition (S1);
\item temperature-dependent elastic Inconel~718 properties.
\end{itemize}
Before any nonlinear run, the LS-DYNA model was verified against Mechanical. The static LC2 state gives a peak von Mises stress of
605.160~MPa against 605.161~MPa at the same location. The LS-DYNA linear buckling factor is 1.1095, +0.13\,\% from the Mechanical value, or
$-$0.64\,\% after the documented large-deformation effect is removed. The mode shape agrees with the Mechanical mode~1 to 0.9999998.

In the nonlinear analysis the temperature rise above 300~K is scaled by a load parameter $\lambda$ up to 1.3. The load steps are fine
($\Delta\lambda$ = 0.005) between $\lambda$ = 1.0 and 1.2, around the expected instability. The imperfection is the normalised first buckling
mode, applied with peak lateral amplitudes of 0.1~mm (case C1), 0.6~mm (C3) and 1.2~mm (C5). These amplitudes are numerical sensitivity
values, not manufacturing tolerances. A perfect-geometry run (C0) serves as a reference path. All four runs terminated normally.

\subsection{Results}
Table~\ref{tab:lsdyna} summarises the three-point imperfection sensitivity study, and Figure~\ref{fig:lsdpath} shows the load paths.

\begin{itemize}
\item \textbf{Load path.} All imperfect cases deflect laterally from the start of loading. The axial compression $N$ (the end reaction)
      flattens towards the critical load as the deflection grows.
\item \textbf{Effect of amplitude below the critical load.} The imperfection amplitude governs where the response leaves the perfect path,
      defined here as $N$ falling 1\,\% below the perfect-geometry value: at $\lambda$ = 1.070 (C1), 0.825 (C3) and 0.375 (C5). It also governs
      the response at the LC2 temperature field ($\lambda$ = 1):
      \begin{itemize}
      \item $N$ is 552.2, 531.2 and 503.5~kN;
      \item the peak von Mises stress is 690, 966 and 1104~MPa;
      \item the end faces translate sideways by 0.85, 3.71 and 5.32~mm.
      \end{itemize}
\item \textbf{Maximum load.} The smallest imperfection reaches a maximum of 601.0~kN at $\lambda$ = 1.185 and then decreases slowly while the
      sway continues to grow. C3 and C5 are still rising at the end of the analysis ($\lambda$ = 1.3), at 572.4 and 545.3~kN.
\item \textbf{Characteristic load.} Southwell plots of the pre-critical branches give 612.3 (C1), 611.1 (C3) and 607.7~kN (C5).
\end{itemize}

%% Section 13 - additional LS-DYNA analysis; values from 15_LS_DYNA_Extension/12C_Nonlinear_Buckling/IMPERFECTION_SENSITIVITY.csv
\begin{table}[htbp]
\centering
\footnotesize
\setlength{\tabcolsep}{3.4pt}
\caption{Additional LS-DYNA nonlinear analysis: three-point imperfection sensitivity under S1 / LC2 (elastic, geometrically nonlinear)}
\label{tab:lsdyna}
\begin{tabular}{l r r r r r >{\raggedright\arraybackslash}p{0.155\linewidth} r r}
\toprule
Case & $A$ [mm] & 1\,\% drop $\lambda$ & $N$ ($\lambda$ = 1) [kN] & VM ($\lambda$ = 1) [MPa] & Yield $\lambda$ & Maximum $N$ ($\lambda \le$ 1.3) & Southwell [kN] & vs $P_{cr}$ \\
\midrule
C1 & 0.1 & 1.070 & 552.2 & 690 & 1.109 & 601.0~kN at $\lambda$ = 1.185 (maximum) & 612.3 & +0.7\,\% \\
C3 & 0.6 & 0.825 & 531.2 & 966 & 1.031 & 572.4~kN at $\lambda$ = 1.3, still rising & 611.1 & +0.5\,\% \\
C5 & 1.2 & 0.375 & 503.5 & 1104 & 0.974 & 545.3~kN at $\lambda$ = 1.3, still rising & 607.7 & $-$0.1\,\% \\
\midrule
C0 (perfect) & 0 & -- & 553.2 & 605.2 & not reached & no instability flagged (720.7~kN at $\lambda$ = 1.3) & \multicolumn{2}{l}{not identifiable} \\
\bottomrule
\end{tabular}
\par\smallskip\parbox{\linewidth}{\footnotesize\raggedright Source: \texttt{15\_LS\_DYNA\_Extension/12C\_Nonlinear\_Buckling/IMPERFECTION\_SENSITIVITY.csv} and \texttt{FINAL\_12C\_RESULTS.md} (additional LS-DYNA analysis); Mechanical reference $\lambda_1$ = 1.108 and $P_{cr}$ = 608.25~kN from \texttt{MASTER\_PROJECT\_DATA.csv} rows M098--M099. $A$: imperfection amplitude. 1\,\% drop: $\lambda$ at which $N$ falls 1\,\% below the perfect path. $N$: axial compression (end reaction); VM: peak von Mises stress. Yield $\lambda$: first local yield, i.e.\ the peak von Mises stress equals the local yield estimate $S_y(T)$ of the elastic model. The perfect-geometry run C0 is a reference path; its $N$ at $\lambda$ = 1 is 0.79\,\% above the Mechanical 548.94~kN because of the large-deformation formulation. Load values are not factors of safety.}
\end{table}


\begin{figure}[htbp]
\centering
\includegraphics[width=\linewidth,keepaspectratio]{LSD_load_lateral_nt.png}
\caption[Additional LS-DYNA analysis: axial compression against lateral end sway]{Additional LS-DYNA analysis: axial compression against the
total relative sway of the end faces for the three imperfection amplitudes and the perfect-geometry reference (left), and the same paths
normalised by the Mechanical critical load and the initial sway (right).}
\label{fig:lsdpath}
\src{Plotted with matplotlib from the LS-DYNA solver output (binout and nodout) of the additional analysis; title strip cropped.}
\end{figure}

\subsection{Comparison with the Mechanical Linear Buckling Result}
The Southwell characteristic-load estimates of 607.7--612.3~kN lie within $-$0.1\,\% to +0.7\,\% of the Mechanical critical load of
608.25~kN ($\lambda_1$ = 1.108), as shown in Figure~\ref{fig:lsdcomp}.

\begin{itemize}
\item \textbf{Imperfection-sensitive response.} The amplitude strongly affects:
      \begin{itemize}
      \item the onset of lateral deformation;
      \item the displacement and stress at a given load;
      \item the maximum load reached within $\lambda \le$ 1.3: 98.8, 94.1 and 89.7\,\% of $P_{cr}$.
      \end{itemize}
\item \textbf{Insensitive characteristic load.} The characteristic global buckling load is relatively insensitive to the tested amplitudes
      and supports the magnitude of the linear eigenvalue prediction.
\item \textbf{Perfect-geometry run.} C0 stays straight beyond $P_{cr}$ (720.7~kN at $\lambda$ = 1.3) without flagging a bifurcation, because
      the solver ignores negative eigenvalues by default. A perfect-geometry bifurcation load is therefore not identifiable from the LS-DYNA
      analysis, and the Mechanical eigenvalue remains the linear perfect-geometry reference.
\item \textbf{Global mode.} In every case the response is a global sideways (guided-sway) bending of the whole duct, the same shape as
      Mechanical mode~1 (Figure~\ref{fig:lsdshape}):
      \begin{itemize}
      \item the mid-span lateral displacement stays below 0.2\,\% of the end-face translation;
      \item the ovalisation is at most 0.041~mm;
      \item there is no local or shell-type mode;
      \item the highest stress is at the outer edge of the inlet face, the critical location of Chapter~\ref{ch:struct}.
      \end{itemize}
\end{itemize}
None of these load values is a factor of safety or a real-world capacity, and the analysis is not an experimental validation.

\begin{figure}[htbp]
\centering
\begin{subfigure}[t]{\linewidth}\centering\includegraphics[width=0.78\linewidth,keepaspectratio]{LSD_mechanical_vs_nonlinear_nt.png}\caption{Mechanical and LS-DYNA linear critical loads; maximum attained $N$ (bars) and Southwell estimates (markers) of the nonlinear cases}\end{subfigure}
\par\medskip
\begin{subfigure}[t]{\linewidth}\centering\includegraphics[width=0.62\linewidth,keepaspectratio]{LSD_amplitude_vs_characteristic_load_nt.png}\caption{Imperfection amplitude against maximum attained load, Southwell estimate and onset load}\end{subfigure}
\caption[Mechanical linear buckling against the additional LS-DYNA nonlinear analysis]{Mechanical linear buckling against the additional
LS-DYNA nonlinear analysis. In (a), 8A denotes the Mechanical linear buckling result of Chapter~\ref{ch:buckling} and 12B G6 the LS-DYNA linear
buckling verification run. Load values are not factors of safety.}
\label{fig:lsdcomp}
\src{Plotted with matplotlib from the LS-DYNA solver output of the additional analysis and the Mechanical result (\texttt{MASTER\_PROJECT\_DATA.csv} M099); title strips cropped.}
\end{figure}

\begin{figure}[htbp]
\centering
\includegraphics[width=\linewidth,keepaspectratio]{LSD_deformed_shape_C5_nt.png}
\caption[Additional LS-DYNA analysis: deformed shape of the 1.2~mm case]{Additional LS-DYNA analysis, case C5 (1.2~mm): lateral deflection
along the duct for $\lambda$ = 1.15--1.30 (left), and the deformed outline at $\lambda$ = 1.30 against the imperfect geometry (right, lateral
axis stretched). Global guided sway, with the ends moving in opposite directions.}
\label{fig:lsdshape}
\src{Plotted with matplotlib from the LS-DYNA d3plot output of the additional analysis; title strip cropped.}
\end{figure}

\subsection{Limitations of the Elastic Model}
The LS-DYNA model uses elastic, temperature-dependent material behaviour. The project contains no defensible temperature-dependent plastic
stress--strain curve for Inconel~718 (task T-036), so the analysis is not elastic--plastic and does not predict post-yield or collapse
behaviour.

\begin{itemize}
\item \textbf{Yield is reached within the load range.} The peak von Mises stress reaches the local yield estimate at $\lambda$ = 1.109
      (C1), 1.031 (C3) and 0.974 (C5). C5 slightly exceeds the yield estimate before $\lambda$ = 1, and all three cases reach yield during
      the extended loading range.
\item \textbf{Consequences.} The later portions of the load paths lie outside the validated elastic constitutive range. The stresses
      computed there are not physical, and no part of the extended response is a collapse load. The Southwell characteristic-load
      comparison, which uses the pre-critical branches, is the meaningful result.
\item \textbf{Other limitations.}
      \begin{itemize}
      \item The imperfection amplitudes are numerical rather than measured.
      \item Only three amplitudes were analysed.
      \item No mesh-refinement or load-step study was made for the nonlinear response.
      \item The end restraint remains the idealised S1 condition discussed above.
      \end{itemize}
\end{itemize}
\FloatBarrier
